%! AP Statistics — Unit 1, Lesson 1.6 — Exit Ticket
\documentclass[10pt]{article}
\usepackage{apstats-article}
\usepackage{apstats-boxes}

%=============================================================================
\begin{document}

\pageheader{Unit 1, Lesson 1.6}{Exit Ticket}
\namedateperiod

\vspace{0.22in}

\begin{tcolorbox}[colback=sky, colframe=navy, boxrule=0.9pt, arc=2mm,
    left=4mm, right=4mm, top=3mm, bottom=3mm]
\small A researcher recorded the \textbf{commute time (in minutes)} for
20 workers. The distribution is roughly \textbf{skewed right}: most
workers commute 10--30 minutes, but a few commute 60--90 minutes.
The approximate center is \textbf{22 minutes} and values range from
\textbf{8 to 88 minutes}. Two workers with commutes over 75 minutes
stand out from the rest of the group.
\end{tcolorbox}

\vspace{0.22in}

\begin{enumerate}[leftmargin=1.5em, itemsep=14pt]

  \item Using the SOCS framework, write a complete description of the
        distribution of commute times for these 20 workers. Address all
        four components in context.\\[6pt]
        \writeline\\[1pt]\writeline\\[1pt]\writeline\\[1pt]\writeline

  \item A classmate says: ``The distribution is skewed right, so the mean
        commute time is probably \textit{less than} the median.'' Is your
        classmate correct? Explain.\\[6pt]
        \writeline\\[1pt]\writeline

\end{enumerate}

\vspace{0.18in}

% ── AP-STYLE ─────────────────────────────────────────────────────────────────
\begin{tcolorbox}[colback=redbg, colframe=redacc, boxrule=0.8pt, arc=2mm,
    left=4mm, right=4mm, top=3mm, bottom=3mm]
\small\textbf{AP-Style Practice}\\[6pt]
A dotplot of the ages (in years) of 15 dogs at an animal shelter shows one
peak near 2--3 years, another smaller peak near 8--9 years, and no obvious
outliers. The ages range from 1 to 10 years.

\vspace{8pt}
Describe the distribution of ages of these 15 dogs using the SOCS framework.
Use appropriate vocabulary and include context in every sentence.\\[8pt]
\writeline\\[1pt]\writeline\\[1pt]\writeline
\end{tcolorbox}

\end{document}
